% New standalone increment; prior increments and GitHub are read-only.
\documentclass[11pt]{article}
\usepackage[T1]{fontenc}\usepackage[utf8]{inputenc}\usepackage{lmodern}
\usepackage[margin=.85in]{geometry}
\usepackage{amsmath,amssymb,amsthm,mathtools,array,longtable,booktabs,microtype,xurl}
\usepackage[hidelinks]{hyperref}
\setlength{\emergencystretch}{3em}\allowdisplaybreaks[2]
\numberwithin{equation}{section}
\newtheorem{theorem}{Theorem}[section]
\newtheorem{lemma}[theorem]{Lemma}
\newtheorem{proposition}[theorem]{Proposition}
\newtheorem{corollary}[theorem]{Corollary}
\newcommand{\E}{\mathbb E}\newcommand{\dd}{\,d}
\newcommand{\KK}{\mathcal K}\newcommand{\BB}{\mathcal B}
\newcommand{\Gp}{G_+}\newcommand{\EE}{\mathfrak E}
\newcommand{\dist}{\mathfrak d}\newcommand{\one}{\mathbf1}
\newcommand{\Kfun}{\mathsf K}\newcommand{\Ffun}{\mathsf F}
\title{Finite-window continuation and initialized error transfer\\in the deep regime}
\author{Standalone analytical increment for independent review}
\date{Version 1 --- 2 November 2026}
\begin{document}\maketitle
\begin{abstract}
Starting at the proved small fixed logistic level, a combined transverse
energy propagates the cross residuals on any fixed time interval.
It supplies bounded strong support, a reached growing weak chart,
primary probe inactivity and the primitive inputs to the probe-sector
energy estimate, without retaining the small-response passage guard.
The exact Gaussian field calculations extend in value to this class.
Their errors, including reaction errors against the current weights,
vanish in the required first-moment norm. Perturbed propagation then
gives initialized error transfer on every fixed centered window,
with the original clock phase and all background populations retained.
Sourcewise estimates are recorded for the next rate-transfer step;
no learned-window rate conclusion or persistence assembly is made here.
\end{abstract}

\section{Setting, sources and the early-window corollary}
\label{an07fwcv1:inventory}
Use normalized time \(t=\tau=\mu_1^2t_{\rm phys}\),
\(\rho\in[13/20,7/10]\), \(\lambda=\rho^2\),
\(Q=\lambda+q^2\), \(S_0=q^{10}\), and \(L=1+\log(1/q)\).
The model is the original balanced untied Gaussian flow,
\(U=\sqrt m\,a_\theta, W=\sqrt m\,a_\psi\), with
\(P_1=N(e_1,q^2I)\), \(P_2=N(\rho e_2,q^2I)\) and mixture halves.
Write \(d\mu=m\,d\lambda_0\) on fixed initial-label cohorts.
Keep the fixed partition
\[
 \KK,\quad G=\Gp\cup\Sigma,\quad\BB,\qquad
 S=\mathcal C_q\subset\Gp,\quad P_q=\Gp\setminus S.
\]
In particular the strong complement is in \(\BB\).
No family is reclassified by its later angle.
Write \(f(X)=\int W(U\cdot X)_+\,d\lambda_0\) and
\(e=X-f\). The source matrices and exact equations are
\[
 R^{(j)}_\theta=\tfrac12\E_{P_j}[eX^T\one_{\{U\cdot X>0\}}],
 \quad R=R^{(1)}+R^{(2)},\quad U'=R^TW,\quad W'=RU.
\]
Thus all source matrices include the mixture half.
At the probe \(\xi_q=(\sin(q\zeta),-\cos(q\zeta))\),
\(E_{\xi_q}=\tfrac12|\xi_q-f(\xi_q)|^2\).
Small \(q\) is chosen also so that \(qR_\zeta\le\pi/2\).
Set \(M=\mu_\KK\), \(N=\mu_S\); these are physical masses,
not the old residual guard or the auxiliary logistic.
Strong and weak coordinates are, respectively,
\[
 (\theta,\psi)=(qx,qy),\qquad
 (\theta,\psi)=(\pi/2-qu,\pi/2+qv).
\]
The weak output convention has the opposite sign.
Let \(n\) be WCS's explicit fixed level, unchanged:
\[
 \begin{gathered}
 n=\tfrac12\min\{h,1/4,
 \lambda_{\min}g_{\min}/[64(980)B_{\max}(20)]\},\\
 \lambda_{\min}=169/400,\quad g_{\min}=1183/41600,
 \quad B_{\max}=e\,2^{1/(2\lambda_{\min})}.
 \end{gathered}
\]
WCS/PSE define its hit \(t_n<t_b\) through the exact auxiliary
logistic starting at \(N_0=q^{1/8}\). Define
\begin{equation}
 \tau_h=t_n+\lambda^{-1}\log\frac{1-n}{n},\qquad
 n_\rho(s)=\frac1{1+e^{-\lambda s}},\qquad
 s_n(\rho)=\lambda^{-1}\log\frac n{1-n}.
 \label{an07fwcv1:time}
\end{equation}
Thus \(t_n=\tau_h+s_n\) exactly. The physical mass \(N\)
will be compared with \(n_\rho\), not declared equal to it.
Fix \(T>\max_\rho s_n\), a bounded nonnegative offset range
\([0,R_\zeta]\), and put
\(D=T-\min_\rho s_n\). All constants below may depend on this
fixed \(D,n,R_\zeta\) and the already fixed entry constants.
They are independent of \(q\). Write
\[
 \chi=1/20,\quad p=603/640<1,\quad
 \gamma_*=289169/924000>3/10,\quad
 \epsilon_{\rm bg}=q^8+L^2q^{\gamma_*}.
\]

\paragraph{Entry data used, with their scopes.}
PSE \texttt{levels,entry}, WCS \texttt{handoff}, CWI
\texttt{entrydistance}, and RB/MP at \(t_n\) give
\begin{align}
 \dist(t_n)&=O_n(q^\chi),& M(t_n)&=1+O(q^{11/10}),&
 N(t_n)&=n+O_n(q^{1/8}),\notag\\
 \sup_\KK(|x|+|y|)&\le C_n,&
 \sup_S(|u|+|v|)&\le C_nq^{-p},& J_G(t_n)&\le C_nq^2,\notag\\
 \mu_{P_q}(t_n)&\le C_nq^{1/4},&
 \mu_\Sigma(t_n)&\le C_nq^{\gamma_*},&
 \mu_\BB(t_n)&\le q^{11/10},\notag\\
 \sup_b|\theta-\psi|(t_n)&\le CqL,&
 \EE_{\Sigma\cup\BB}(t_n)&\le Cq^2\epsilon_{\rm bg}.
 \label{an07fwcv1:entry}
\end{align}
Here \(J_A=\tfrac12\int_A(U_1^2+W_1^2)\) and
\(E_A=\tfrac12\int_A(U_2^2+W_2^2)\), so
\(E_A+J_A=\mu_A\). In addition, PSE \texttt{inactive}'s proof,
using RB/RCET only before this entry, supplies
\begin{equation}
 U_2,W_2>0\text{ on }\Gp,\qquad
 \frac{q|U_1|}{U_2}\le Cq^{597/640}\quad\text{at }t_n.
 \label{an07fwcv1:primaryentry}
\end{equation}
Its regularized endpoint labels, including \(b=\pi\), are included.
The distance in \eqref{an07fwcv1:entry} uses physical current measures;
the fixed incoming-law \(L^2\) estimate and the proved weight conversion
are used only to obtain this entry value.

\small
\begin{longtable}{>{\raggedright\arraybackslash}p{.22\textwidth}
 >{\raggedright\arraybackslash}p{.41\textwidth}
 >{\raggedright\arraybackslash}p{.28\textwidth}}
\toprule Source & All hypotheses used here & Interval and matching\\\midrule\endhead
PSE \texttt{levels,entry}
& Original initialized deep regime; fixed WCS level; compact level set
 inside \((0,n]\); same auxiliary logistic; small \(q\).
 Rates have WIT's positive offset range and selector convention.
& Only at \(t_{n'}\le t_n<t_b\). Supplies entry and the early
 corollary, not later evolution.\\
WCS \texttt{handoff}; RB/MP/PSE entry budgets
& Their completed initialized regime, unchanged constant order,
 positive radial measures, reached charts and \(t\le t_b\).
& Used only at \(t_n\), including current-weight conversion,
 physical masses, \(J_G\), outside mass and alignment.\\
P exact flow and loss
& Original Gaussian model, balanced positive initialized radii,
 finite label law, normalized time, finite moments.
& Exact identities on each finite continuation. Loss gives bounded
 residual matrices and existence; no sign of \(c\) is assumed.\\
CFA/SRR core Gaussian calculations
& Bounded strong donor/receiver chart; bounded mass and moments;
 primary positive \(P_2\) margins; remaining donors charged in
 absolute value. Capping and \((K_O)_+\) are not used here.
& Their passage theorems are not extended. The value calculations
 are restated with physical \(M\), signed \(K_O\), and the new
 energy bounds in Section 4.\\
WCS \texttt{IBP, P2, overlap, reduction}
& Reached chart \(q(|u|+|v|)=o(1)\); positive current selected
 measure; bounded donor second moments, core chart, transverse
 energy, omitted-energy and outside charges.
& Exact identities re-used; the originally nonnegative-input
 Gaussian estimates are proved for both signs below.
 No use of \(u\ge0\), \(Y\ge0\), \(c\le c_b\) after entry.\\
CWI \texttt{propagation}
& Defined leading velocities/reactions; positive measures,
 masses \(\le2\); bounded strong support and comparator;
 \(R\ge2\), \(0<\rho\le1\); finite weak first moments.
& Not applied directly to the original flow. Its one-sided
 comparison is repeated with velocity and weighted reaction
 errors, supplied in Sections 4--5.\\
PSE \texttt{conditional}; CWI \texttt{split,ratecharge}
& Bounded source matrices, cross entries \(O(q)\),
 \(R^{(1)}_{22}=O(q^2)\), lower Gaussian gate, alignment
 \(O(qL)\); fixed cohorts, bounded nonnegative offsets.
 Split also needs inactive \(\Gp\), \(f_\KK=O(q)\), and
 \(L^2\mu_{\Sigma\cup\BB}\to0\).
& All re-proved on the new window. No derivative of a receiver
 gate is used. Rate statements retain the selector/a.e. convention.\\
CWI static part of \texttt{rate}
& Positive strong measure, \(M\le2\), bounded strong chart
 and comparator, bounded offset.
& Its energy proof is restated with \((\zeta-X)_+\), so no
 probe gap is required for the error observable.\\
H2 A-error-transfer; LP \texttt{compact-transfer}
& Actual initialized flow, all populations, same analytically
 matched center; each fixed finite centered window; uniform
 compact parameters and offsets; original times nonnegative.
& Proved below. LP uses \(t_h^{\rm phys}=\tau_h/\mu_1^2\).
 Neither clearing nor the diagonal argument is run here.\\
WIT rate/window consumers
& In addition: \(\mathcal P_*\), learned windows, bounded
 comparator, positive gap \(g_0\), sourcewise precision.
& Handoff only; no A-rate-transfer conclusion is asserted.\\
\bottomrule
\end{longtable}\normalsize

\begin{corollary}[Already proved early centered windows]
\label{an07fwcv1:early}
For every fixed compact \(K\subset(-\infty,\min_\rho s_n]\),
the full initialized A-error-transfer holds uniformly for
\(s\in K\), the deep ratio box, and \(0\le\zeta\le R_\zeta\).
On WIT's offset range PSE also supplies the entry source-rate
comparison on these windows, in its stated rate convention.
\end{corollary}
\begin{proof}
Set \(n'=n_\rho(s)\). Compactness and \(\lambda\in[169/400,49/100]\)
give \(0<n_-\le n'\le n\), uniformly. PSE \texttt{levels,entry}
applies on this single compact level set, and
\(t_{n'}=\tau_h+\lambda^{-1}\log(n'/(1-n'))=\tau_h+s\)
exactly. This proves the stated conclusion on WIT's offsets.
For arbitrary bounded nonnegative offsets, use the same entry
first-moment estimate, the core static positive-part expansion,
and PSE's background bound, both uniform on \([0,R_\zeta]\).
The static Lipschitz argument is displayed in Section 6 and needs
no gap. Thus the error conclusion has the larger offset range;
no corresponding gap-free rate assertion is inferred.
Since \(\tau_h\to\infty\), these original times are positive.
\end{proof}

\section{What is replaced after the passage}
\label{an07fwcv1:rangeaudit}
Neither the LCRC guards nor the RB clock are invoked after \(t_n\).
Their functions in the old reductions and their replacements are:
\begin{itemize}
\item The small \(E_g\) and \(c_b\) were essential to the long
 passage closure and CLOCK attraction. They are not hypotheses of
 the Gaussian value identities. Bounded total mass and the new
 transverse estimate replace their value-budget roles.
\item The \(P_2\) primary-gate tail previously came from RCET's
 growing margin. Its entry margin is propagated for a fixed duration
 directly by the exact angular equations; no later residual clock is used.
\item The old secondary bound \(q^5+q^{\Gamma_\nu}H\) and all strip
 gain theorems are used only at entry. Fixed-label radial comparison
 propagates their small masses; the combined transverse energy
 pays their possibly unbounded per-mass fields.
\item The capped core measure \(\widetilde M\le1\) and the positive
 part \((K_O)_+\) served CLOCK's attraction proof. Here the reduction
 retains physical \(M\) and signed \(K_O\). No overshoot cap or
 positive strong output moment is needed for finite-time comparison.
\item WCS's return theorem and local contraction cannot be extended
 to arbitrary later masses. Instead the reached weak chart is
 propagated, Gaussian estimates are extended to signed \(u\),
 and CWI's one-sided comparison is used. There is no claim that
 the actual weak input stays nonnegative.
\item Small integrated passage budgets are not pointwise errors.
 The remainders below come from the displayed value identities
 and the newly proved instantaneous bounds.
\end{itemize}

\section{Primitive finite-time continuation}
\label{an07fwcv1:continuation}
Let \(I_q=[t_n,\tau_h+T]\), of length at most \(D\).
All supremum bounds in this section are on this interval.

\begin{lemma}[Loss and fixed-label masses]
\label{an07fwcv1:massrough}
The exact original flow exists on \(I_q\). For \(q\le1/2\),
one may use the model bound \(\|R_\theta\|,\|R^{(j)}_\theta\|\le4\),
and, for every fixed cohort \(A\),
\begin{equation}
 e^{-8D}m(t_n,b)\le m(t,b)\le e^{8D}m(t_n,b),\qquad
 \mu_A(t)\le e^{8D}\mu_A(t_n),\qquad
 S_{\rm tot}(t)\le\bar S:=3e^{8D}.
 \label{an07fwcv1:roughmass}
\end{equation}
\end{lemma}
\begin{proof}
Loss monotonicity gives \(\|X-f\|_{P_j}\le2\), as in PSE.
Cauchy--Schwarz with the bounded Gaussian second moments bounds
the gated matrices, source by source. The exact radial identity
\((\log m)'=2a_\psi^TRa_\theta\) gives the displayed factors.
Entry total mass is at most three for small \(q\).
At fixed positive \(q\), the Gaussian field is locally Lipschitz
in angles and in the bounded-mass characteristic representation;
the squared radii stay positive and bounded on finite intervals.
These bounds preclude a finite-time characteristic escape and
extend the local flow through \(I_q\). This is an existence
argument, not the desired scaled-chart estimate.
\end{proof}

\begin{lemma}[Combined transverse energy closes the cross residuals]
\label{an07fwcv1:transverse}
Define
\[
 \mathcal Z=E_\KK+J_G
 =\tfrac12\int_\KK(U_2^2+W_2^2)
  +\tfrac12\int_G(U_1^2+W_1^2),\qquad
 \mathcal X=\|e_2\|_{P_1}+\|e_1\|_{P_2},\quad e=X-f.
\]
Then on \(I_q\)
\begin{equation}
 \mathcal Z\le C_Dq^2,\qquad
 \mathcal X\le C_Dq,
 \quad |R_{12}^{(j)}|+|R_{21}^{(j)}|\le C_Dq,
 \quad |R_{22}^{(1)}|\le C_Dq^2.
 \label{an07fwcv1:primitive}
\end{equation}
No chart, primary gate sign, small response or small weak energy
is assumed in proving these estimates.
\end{lemma}
\begin{proof}
For any cohort \(A\), balance and Cauchy--Schwarz give
\begin{align}
 |f_{A,1}(X)|&\le2J_A|X_1|+2\sqrt{E_AJ_A}|X_2|,\notag\\
 |f_{A,2}(X)|&\le2\sqrt{E_AJ_A}|X_1|+2E_A|X_2|.
 \label{an07fwcv1:outputenergy}
\end{align}
Use the second bound for the core on \(P_1\) and the first
on \(P_2\), and the reverse roles of the small energies for
\(G\). The target cross coordinates have norm \(q\).
The arbitrary \(\BB\) output has norm at most \(C\mu_\BB\).
Hence, with a model constant,
\begin{equation}
 \mathcal X\le C[q(1+\bar S)+\sqrt{\bar S\mathcal Z}+\mu_\BB].
 \label{an07fwcv1:crossfromenergy}
\end{equation}
At every receiver, bounded diagonal residuals and these two
cross norms imply \(|R_{12}|+|R_{21}|\le C(q+\mathcal X)\),
also sourcewise. The exact Cartesian equations give
\[
 (E_\KK)'=\int_\KK
 [2R_{22}U_2W_2+R_{12}U_2W_1+R_{21}W_2U_1],
\]
and the corresponding formula for \((J_G)'\) has
\(2R_{11}U_1W_1\) and the same cross products on \(G\).
Their matrices are evaluated at their own receivers.
Consequently
\[
 D^+\mathcal Z\le C\mathcal Z
     +C(q+\mathcal X)\sqrt{\bar S\mathcal Z}
 \le C_{\bar S}[\mathcal Z+q^2+\mu_\BB^2].
\]
Entry has \(\mathcal Z(t_n)=O(q^2)\) by its core chart
and \(J_G\) bound. Lemma \ref{an07fwcv1:massrough} gives
\(\mu_\BB^2\le e^{16D}q^{11/5}\le e^{16D}q^2\).
Gronwall proves \(\mathcal Z\le C_Dq^2\), then
\eqref{an07fwcv1:crossfromenergy} gives \(\mathcal X\le C_Dq\).
Finally \(|R_{22}^{(1)}|\le\frac12q\|e_2\|_{P_1}\).
All coefficients were fixed before any proposed scaled-energy guard.
\end{proof}

\begin{lemma}[Charts, alignment, primary gates and moments]
\label{an07fwcv1:charts}
There are fixed constants such that throughout \(I_q\)
\begin{gather}
 \sup_b|\theta-\psi|\le C_DqL,\qquad
 \sup_\KK(|x|+|y|)\le R_A,\qquad
 \sup_S(1+|u|+|v|)\le A_Dq^{-p},\label{an07fwcv1:chartsizes}\\
 \int_S(1+u^2+v^2)\dd\mu\le C_D,\qquad
 \sin\theta,\sin\psi\ge c_Dq^{43/640}\quad\text{on }\Gp.
 \label{an07fwcv1:momentsgates}
\end{gather}
Every primary label is probe-inactive for \(0\le\zeta\le R_\zeta\).
Selected labels satisfy \(q|u+\zeta|<\pi/2\); their chart is
reached, but no post-entry sign of \(u\) is asserted.
\end{lemma}
\begin{proof}
The exact difference equation is
\(\Delta'=-\operatorname{tr}R\sin\Delta+(R_{12}-R_{21})\cos\Delta\).
Bounded matrices and \eqref{an07fwcv1:primitive} give
\(D^+|\Delta|\le C|\Delta|+C_Dq\) while \(|\Delta|<1\).
From the entry value \(O(qL)\), fixed-time Gronwall strictly
closes that auxiliary guard. It also gives
\(|W_2-U_2|\le C_DqL\sqrt m\).

Near the strong axis the exact angular equations give
\(D^+(|\theta|+|\psi|)\le C(|\theta|+|\psi|)+C_Dq\)
while both angles have magnitude below \(1/4\).
Starting at \(O(q)\), they stay within \(R_Aq\) for a fixed
\(R_A\), strictly inside that physical chart for small \(q\).
Near the positive weak axis the same calculation with the
distances to \(\pi/2\) gives
\(D^+(|\theta-\pi/2|+|\psi-\pi/2|)
\le C(|\theta-\pi/2|+|\psi-\pi/2|)+C_Dq\).
Entry is \(O(q^{1-p})\), which proves the selected chart.
On it, \(|qu|,|qv|\le1/2\) and \(|\sin z|\ge|z|/2\),
so its second moment is bounded by \(C J_S/q^2\le C_D\).

For the whole primary family, \eqref{an07fwcv1:primaryentry}
gives \(\sin\theta(t_n)\ge c q^{43/640}\).
The alignment bound gives the same for \(\psi\).
As long as both are positive, direct differentiation of
\(\sin\theta\), using the bounded diagonals and cross entries,
gives
\[
 (\sin\theta)'\ge-C\sin\theta-C_DqL,
\]
and likewise for \(\psi\). For example replace \(\sin\psi\)
by \(\sin\theta+O(qL)\) in the diagonal terms; the remaining
terms cost \(C_Dq\). Thus both stay above
\(e^{-CD}c q^{43/640}-C_DqL\ge c_Dq^{43/640}>0\).
This strictly prevents their first zero, including boundary labels.
It follows that \(q|\cot\theta|\le C_Dq^{597/640}\), so
\(a_\theta\cdot\xi_q<0\) uniformly on bounded nonnegative offsets.
For \(S\) the stronger reached-chart bound gives the asserted
\(q|u+\zeta|\) inequality. All classifications use current angles.
\end{proof}

\begin{lemma}[Physical mass and response equations]
\label{an07fwcv1:massfine}
Uniformly on \(I_q\),
\begin{align}
 g_\KK(t,b):=(\log m)'&=1-M+\epsilon_\KK,
 &|\epsilon_\KK|&\le C_D(q^2+\mu_\BB),\notag\\
 c&=N+O_D(q^{1/4}),\label{an07fwcv1:cN}\\
 g_S(t,b)&=\lambda(1-N)+\epsilon_S,
 &|\epsilon_S|&\le C_D[q^{1/4}+q^2(1+|u|+|v|)^2].
 \label{an07fwcv1:reactions}
\end{align}
In particular
\begin{equation}
 \sup|M-1|\le C_Dq^{11/10},\qquad
 \sup|N-n_\rho(t-\tau_h)|\le C_D(q^{1/8}+q^{1/4}),
 \quad M,N<3/2
 \label{an07fwcv1:masslogistic}
\end{equation}
for small \(q\). The physical logit has the same \(o(1)\) phase
on this fixed window.
\end{lemma}
\begin{proof}
On \(P_1\), all core gates have exponentially small deficits,
with \(\cos\theta\ge1/2\). Their longitudinal donor coefficient
is \(M+O_D(q^2)\). Equation \eqref{an07fwcv1:outputenergy}
gives \(\|(f_G)_1\|_{P_1}=O_D(q^2)\);
the arbitrary outside contribution costs \(C\mu_\BB\).
The core radial \(P_1\) contraction is therefore
\(1-M+O_D(q^2+\mu_\BB)\). On \(P_2\), core input is
\(O_D(q)(|Z_1|+|X_2|)\); the first residual has norm \(O_D(q)\)
and the core output second component is \(O_D(q)\).
The two radial contractions cost \(O_D(q^2)\). This proves
the first line, with its mixture normalization.

For the global response \(c=Q^{-1}\E_{P_2}[f_2X_2]\),
the selected full-gate expression is
\(\int_S W_2U_2=N+O(J_S)\), with exponentially small error.
Core response is \(O_D(q^2)\) by conditioning on
\(X_2=\rho+qZ_2\) in the compact core moment.
For any other cohort the exact absolute response charge gives
\(|c_A|\le E_A+2q\sqrt{E_AJ_A/Q}\le C\mu_A\).
The mass comparison gives \(\mu_{P_q}=O_D(q^{1/4})\),
\(\mu_\Sigma=O_D(q^{\gamma_*})\), \(\mu_\BB=O_D(q^{11/10})\).
This proves \eqref{an07fwcv1:cN} without assuming a response sign.

For selected receivers the \(P_2\) gate is exponentially full.
The exact full-gate moment is \(B_{22}=Q(1-c)\),
\(B_{11}=O_D(q^2)\), \(B_{12}=O_D(q)\),
\(B_{21}=O_D(q^3)\). For the last entry, Stein gives
\(-q^2\int W_2U_1\Pr\{U\cdot X>0\}\), and core,
\(G\), outside charges are \(O_D(q),O_D(q),O_D(q^{11/10})\).
The radial contraction thus equals \(Q(1-c)\) plus
\(O_D(q^2(1+|u|+|v|)^2)\). On \(P_1\), use
\(R_{22}^{(1)}=O_D(q^2)\), cross entries \(O_D(q)\),
and horizontal receiver components \(O(q|u|),O(q|v|)\)
to get that same error size. Replace \(Q,c\) by \(\lambda,N\)
using \eqref{an07fwcv1:cN}; the response is bounded by the
rough total mass, so this replacement has the stated cost.

Integrating the radial laws against current weights gives
\(M'=M(1-M)+O_D(q^{11/10})\) and
\(N'=\lambda N(1-N)+O_D(q^{1/4}+q^2)\), using the second
moment, not the weak supremum. Both masses were already bounded
by \(\bar S\), so finite-time scalar stability gives
\eqref{an07fwcv1:masslogistic} from \eqref{an07fwcv1:entry}.
In particular masses \(\le2\) are now conclusions, not a
circular hypothesis of the field comparison. On a fixed window
\(n_\rho\) stays uniformly away from zero and one; hence the
same estimate puts the actual mass strictly between zero and one
for sufficiently small \(q\) and controls its logit difference. No derivative of
an entry phase error was taken.
\end{proof}

\paragraph{First-exit implementation and constant order.}
The estimates above also give a single strict-guard proof if desired.
First fix the inherited WPT/RB--WCS constants and \(n\), then
\(T,R_\zeta,D,\bar S\). Fix an energy bound twice the explicit
Gronwall majorant from \eqref{an07fwcv1:crossfromenergy}; next
twice the resulting cross-residual constant, then twice the
alignment and angular-comparison bounds. Use guards
\[
 \mathcal Z<2C_Zq^2,\quad \mathcal X<2C_Xq,\quad
 |\Delta|<2C_\Delta qL,\quad
 |x|+|y|<2R_A,\quad 1+|u|+|v|<2A_Dq^{-p},
 \quad M,N<2,
\]
and primary lower guards one half of the proved positive lower
bound. Fixed-cohort masses have guards twice
\(e^{8D}\mu_A(t_n)\); the rough total mass can use
\(2\bar S\). Every entry is strict. The energy majorant depends
only on the independent rough mass bound, not on the energy or
cross-residual guard constants. Subsequent bounds depend only
on their previously fixed predecessors; the mass improvement
uses the exact radial laws, not the first-moment comparison.
Every guard is improved by its displayed bound. Supremum guards
are continuous (compact label sets, or their closures with the same
endpoint estimates); the integrals are continuous by domination.
At a first exit the preceding estimates give a contradiction.
Choose \(q\) last to make all physical-chart, lower-margin and
mass inequalities strict. No \(c_b\) or small-energy guard is
continued.

\section{Original fields and reactions in the first-moment norm}
\label{an07fwcv1:reduction}
Use the physical core and selected measures, and define
\[
 Y=\int_\KK y\dd\mu,\quad V=\int_S v\dd\mu,\quad
 K_s=\int_\KK\Kfun(\rho x)\dd\mu,\quad K_w=\int_S\Kfun(u)\dd\mu,
\]
\[
 \begin{gathered}
 \Kfun(z)=\phi(z)+z\Phi(z),\quad
 \Ffun(z,w)=\phi(\min(z,w))+w\Phi(\min(z,w)),\\
 \mathcal S(x)=\int_\KK\Ffun(\rho x,\rho\widetilde x)\dd\mu,
 \quad\mathcal W(u)=\int_S\Ffun(u,\widetilde u)\dd\mu.
 \end{gathered}
\]
These are positive measures, but their coordinate means need not
have either sign. Their first moments are bounded by Section 3.
For \(O=G\cup\BB\), the exact non-core moments used in the
Gaussian calculation are
\[
 V_O=-\frac{\E_{P_2}[(f_O)_1X_2]}{qQ},\qquad
 K_O=\frac{\E_{P_1}[(f_O)_2X_1]}q.
\]
P's leading weak velocity \(F_w=(F_u,F_v)\), with these
physical core and selected measures, is explicitly
\begin{align*}
 2F_u&=\phi(u)-\mathcal W(u)-[Y+(1-M)v]\Phi(u)
                  -\lambda(1-N)u,\\
 2F_v&=\rho K_s-\lambda V-\lambda(1-N)v-(1-M)\Kfun(u).
\end{align*}

\begin{lemma}[Energy charges and signed Gaussian value calculations]
\label{an07fwcv1:charges}
The omitted \(P_q,\Sigma\) moments and selected weak functionals
cost, uniformly in the receiver,
\begin{equation}
 C_D(\mu_A+q^{-1}\sqrt{\mu_AJ_A})
 \le C_D(\mu_A+\sqrt{\mu_A}),\qquad A=P_q,\Sigma.
 \label{an07fwcv1:nonselectedcharge}
\end{equation}
They are respectively \(O_D(q^{1/8})\) and
\(O_D(q^{\gamma_*/2})\). The arbitrary \(\BB\) field costs
\(C_D\mu_\BB/q=O_D(q^{1/10})\).
The core source velocities have P's physical leading form,
with value errors bounded by \(C_Dq^{1/10}\), source by source:
\begin{align}
 2x'_1&=-(1-M)x+e_{x1},\notag\\
 2y'_1&=-(1-M)y-Y-K_w+e_{y1},\notag\\
 2x'_2&=\rho\phi(\rho x)-\rho\mathcal S(x)
       +\lambda[V+(1-N)y]\Phi(\rho x)+e_{x2},\notag\\
 2y'_2&=\rho(1-N)\Kfun(\rho x)+e_{y2}.
 \label{an07fwcv1:corefield}
\end{align}
The core source reactions are
\(g_{\KK,1}=1-M+O_D(q^2+\mu_\BB)\),
\(g_{\KK,2}=O_D(q^2)\).
For selected receivers the leading weak velocities \(F_w\)
are P's \texttt{pop-u,pop-v} with these physical moments and
masses. If \(r=1+|u|+|v|\), then
\begin{equation}
 |(u',v')-F_w(u,v)|_1
 \le C_D[q^{1/10}+q r+q^2r^3+q^{1/4}r].
 \label{an07fwcv1:weakerror}
\end{equation}
No non-selected weak field is charged by its raw mass alone.
The \(\BB\) alternative is the exact total-quadratic-energy charge
\(C_D(E_\BB+J_\BB)/q=C_D\mu_\BB/q\), not a bounded
per-mass scaled-field assumption.
\end{lemma}
\begin{proof}
Equation \eqref{an07fwcv1:outputenergy}, inserted into the
\(P_1\) weak functional \(q^{-1}\E[f_{A,2}\one_{\mathcal A_\theta}]\),
where \(\mathcal A_\theta\) is the receiver gate,
and the two regression moments defining \(V_O,K_O\), gives
\eqref{an07fwcv1:nonselectedcharge} by Gaussian moments.
This repeats SRR \texttt{momentcharge} and WCS
\texttt{omittedfunctional} with their actual hypotheses.
It uses no chart or input sign for the omitted donors.
For \(\BB\), \(|f_\BB(X)|\le\mu_\BB|X|\) directly gives
the unscaled matrix cost \(C\mu_\BB\); angular scaling costs
\(1/q\). No outside energy theorem is extended past its endpoint.

Here is the core remainder audit, using physical mass and signed
moments throughout. Remove core \(P_1\) donor and receiver gates
at exponential cost, and keep its four donor coefficients.
Their compact-chart expansions are
\(T_{11}=M+O_D(q^2)\), \(T_{21}=qY+O_D(q^3)\),
\(T_{12}=O_D(q)\), \(T_{22}=O_D(q^2)\).
The exact \(P_1\) Stein calculation in CFA gives the first
line of \eqref{an07fwcv1:corefield}, with errors
\(C_D[J_G+q\sqrt{E_GJ_G}+q^2(1+E_G+|K_O|)]\)
and the outside cost. These are \(O_D(q^2)+O_D(\mu_\BB/q)\).
For \(P_2\), conditioning on \(X_2\) gives the compact core
Gaussian kernels \(\rho\phi-\rho\mathcal S\) and
\(\rho\Kfun\), with \(O_D(q^2)\) error in scaled value.
Primary outputs are linear up to a Gaussian negative-part tail.
Section 3 gives a donor margin at least \(c_Dq^{-597/640}\);
all polynomial multiples of its tail
\(\exp[-c_Dq^{-597/320}]\) vanish faster than any fixed power.
Subtracting the \(X_2\) regression leaves the linear cross
charges \(J_G+q\sqrt{E_GJ_G}\), not a zero selected moment.
Secondary selected moments cost
\(C_D(E_\Sigma+J_\Sigma+q^{-1}\sqrt{E_\Sigma J_\Sigma})\).
This is \(O_D(q^{\gamma_*/2}+q^2)\) by the proved masses
and transverse energy. The \(\BB\) term costs \(\mu_\BB/q\).
All trigonometric and compact-kernel corrections cost
\(q^2\) times bounded masses/moments. This is precisely the
value calculation behind CFA/SRR, now justified by new bounds;
neither \(c\le c_b\) nor small \(E_G\) enters it. Keep
\(K_O\) signed and do not cap \(M\).

For clarity the selected moment replacements are also small
without weak compact support. On its reached chart, exact Gaussian
regression and the \(P_1\) formula give
\[
 V_S=\int_Sm\frac{\sin(qv)}q\cos(qu)+\text{exponential tail},
\]
\[
 K_S=\int_Sm\cos(qv)
 \left[\Kfun\!\left(\frac{\sin(qu)}q\right)
       +q\sin(qu)\Phi\!\left(\frac{\sin(qu)}q\right)\right].
\]
Thus \(|V_S-V|+|K_S-K_w|\le C_Dq^2\int_S r^3\dd\mu
 +\text{exponential tail}=O_D(q^{2-p})\).
The full non-core moments differ from these by
\eqref{an07fwcv1:nonselectedcharge} and \(\mu_\BB/q\).
Replacing \(c\) by \(N\) costs \(O_D(q^{1/4})\) on the core
chart, by \eqref{an07fwcv1:cN}. Every exponent just listed is
at least \(1/10\). The core radial calculation was given in
Lemma \ref{an07fwcv1:massfine}, source by source.

For the weak receivers, the WCS tangent integration by parts is
exact for both signs. With \(s_u=\sin(qu),c_u=\cos(qu)>1/2\),
rotate to independent normals
\(Z_\parallel=s_uZ_1+c_uZ_2\),
\(Z_\perp=c_uZ_1-s_uZ_2\). The receiver gate depends only
on \(Z_\parallel\); integrate by parts only in \(Z_\perp\).
This gives WCS \texttt{IBP}'s four moments with
\(r_u=s_u/q\), and its exact donor derivative identity.
For signed \(u\), \(r_u\) has the same sign as \(u\) and
\(|u|/2\le|r_u|\le|u|\). Consequently
\begin{equation}
 |\Phi(r_u)-\Phi(u)|+|\phi(r_u)-\phi(u)|\le Cq^2,
 \quad |\Kfun(r_u)-\Kfun(u)|\le Cq^2|u|^3.
 \label{an07fwcv1:signedgaussian}
\end{equation}
Indeed \(|r_u-u|\le q^2|u|^3/6\) and the intervening
Gaussian factors are bounded by a polynomial times
\(e^{-u^2/8}\). This handles the negative side as well as
the positive side. In particular no quartic
\(q^2|v||u|^3\) is generated by a crude bound on \(\Phi\).

In the weak overlap calculation replace a donor's positive part
by \((Z_2+\widetilde u)_+\). Its integrated error, and the
output cosine error, cost
\(C_D[q\int r+q^2\int r^3]=O_D(q+q^{2-p})\).
Conditioning on \(Z_1\), changing the receiver threshold costs
\(C_D(q|u|+q^2|u|^3)\): the integrand
\((z+\widetilde u)_+\phi(z)\) is bounded by
\(C(1+|\widetilde u|)\) for every signed donor and receiver.
The exact remaining overlap is \(\mathcal W(u)\).
All donor first and second moments used here are bounded by
\eqref{an07fwcv1:momentsgates}; cubic moments cost only its
proved supremum times the second moment.

The core \(P_1\) coefficients then give
\(\phi(u)-Y\Phi(u)-(1-M)v\Phi(u)\) and
\(-(1-M)\Kfun(u)\). Errors have degree at most three with
coefficient \(q^2\), by \eqref{an07fwcv1:signedgaussian}.
The selected weak functional gives \(-\mathcal W(u)\);
the omitted functionals have already been charged.
On \(P_2\), the receiver gate is exponentially full, and
the exact moment matrix used in Lemma \ref{an07fwcv1:massfine}
gives \(-\lambda(1-c)u\) and
\(\rho K_s-\lambda V_O-\lambda(1-c)v\),
with \(C_Dq^2r^3\) error. Replacing \(V_O\) and \(c\)
by \(V\) and \(N\) gives exactly
\eqref{an07fwcv1:weakerror}. All original residual producers
are retained in these charges. The calculation never uses
the old positive-input return, small \(E_g\), or local strong
attraction.
\end{proof}

\begin{proposition}[Integrated value and weighted reaction errors]
\label{an07fwcv1:errors}
Let \(F_s,F_w\) be P's leading velocities with physical
\(M,N\) and current measures. Put
\(\varepsilon_s=(x',y')-F_s\),
\(\varepsilon_w=(u',v')-F_w\), and use the reaction defects
\(\epsilon_\KK,\epsilon_S\) from \eqref{an07fwcv1:reactions}.
For any fixed bound on comparison coordinates, there is a fixed
constant such that
\begin{align}
 \mathcal R(t):={}&\int_\KK
    [|\varepsilon_s|_1+(1+|x-X|+|y-Y_*|)|\epsilon_\KK|]\dd\mu
       \notag\\
 &+\int_S[|\varepsilon_w|_1+
          (1+|u-U_*|+|v-V_*|)|\epsilon_S|]\dd\mu
 \le{}&C_Dq^{1/10},\qquad
 \int_{I_q}\mathcal R\le C_DDq^{1/10}.
 \label{an07fwcv1:reductionnorm}
\end{align}
\end{proposition}
\begin{proof}
With \(r=1+|u|+|v|\), the reached chart and transverse energy give
\[
 \int_Sr^2\dd\mu\le C_D,\qquad
 \int_Sr^3\dd\mu\le C_Dq^{-p},\qquad
 q\int_Sr+q^2\int_Sr^3\le C_D(q+q^{2-p}).
\]
The additional weak response term costs
\(q^{1/4}\int_Sr\le C_Dq^{1/4}\).
For reaction-weighted distance, use
\(1+|u-U_*|+|v-V_*|\le C_Dr\) and
\eqref{an07fwcv1:reactions}, giving
\(C_D(q^{1/4}\int r+q^2\int r^3)\).
The core terms are bounded pointwise on its fixed chart.
Since \(2-p=677/640>1\), all these powers dominate
\(q^{1/10}\). This is an integral against current mass;
using the weak supremum cubed would not prove it.
\end{proof}

\section{Perturbed first-moment propagation with current weights}
\label{an07fwcv1:propagationsection}
The selected leading solution has strong mass one and weak mass
\(n_\rho(s)\). Its coordinates are bounded on this fixed window,
uniformly in the ratio box. To verify the needed bound without a
trajectory premise, substitute the coherent atoms into P's system:
\(|(X,Y_*,U_*,V_*)'|_1\le C(1+|(X,Y_*,U_*,V_*)|_1)\),
because \(0<n_\rho<1\), \(|\Phi|\le1\),
\(|\phi|\le\phi(0)\), and \(\Kfun(z)\le\phi(0)+|z|\).
The selected branch exists at \(n\) by P/WCS and has uniformly
bounded entry. Gronwall gives both forward existence and a fixed
coordinate bound on \([s_n,T]\).
Fix \(R\ge2\) larger than this bound, \(R_A\), and \(R_\zeta\).

\begin{theorem}[Original first-moment comparison]
\label{an07fwcv1:propagation}
At physical time \(t=\tau_h+s\), define
\[
 D_s=\int_\KK(|x-X|+|y-Y_*|)\dd\mu,\qquad
 D_w=\int_S(|u-U_*|+|v-V_*|)\dd\mu,
\]
\[
 \dist=|M-1|+|N-n_\rho|+D_s+D_w.
\]
Then
\begin{equation}
 \dist(s)\le e^{64R(s-s_n)}
 \left[\dist(s_n)+\int_{s_n}^s\mathcal R(\tau_h+v)\dd v\right]
 =O_{T,n}(q^{1/20})
 \quad(s_n\le s\le T).
 \label{an07fwcv1:distbound}
\end{equation}
The masses and distances use the original current weights.
In particular the actual selected mass has vanishing mass/logit
phase error on the same logistic clock.
\end{theorem}
\begin{proof}
All structural hypotheses of CWI's leading comparison now hold:
positive measures, \(M,N\le2\), bounded strong support and
comparator, finite weak moments, same \(\rho\). Its strong
velocity comparison is \(8R(d_s+\dist)\).
For completeness the weak one-sided step, valid for all real \(u\),
is
\[
 A'(u)=-\tfrac12[(u+Y)\phi(u)+\mathcal W'(u)
                           +\lambda(1-N)]\le R+1\le2R,
\]
where
\(A(u)=\tfrac12[\phi(u)-\mathcal W(u)-Y\Phi(u)
                         -\lambda(1-N)u]\) and
\(\mathcal W'(u)=\phi(u)\int_S(\widetilde u-u)_+\dd\mu\ge0\).
Use \(|Y|\le2R\), \(|u|\phi(u)\le1\),
\(\lambda|1-N|\le1\), and compare donor moments at bounded
\(U_*\) first. The remaining \((1-M)v\Phi(u)\) and
\((1-M)\Kfun(u)\) terms cost at most the receiver distances
plus \(R|M-1|\), since \(|1-M|\le1\).
Exactly as in CWI this yields
\[
 D^+d_s\le8R(d_s+\dist)+|\varepsilon_s|_1,
 \qquad D^+d_w\le8R(d_w+\dist)+|\varepsilon_w|_1.
\]
The reference donor differences are controlled by physical
mass differences and the two first moments; no weak support
bound or sign of its input is used in this comparison.

Now differentiate the integrals with their actual reactions.
The leading reaction coefficients \(1-M,\lambda(1-N)\)
have absolute value at most one. Their defects add respectively
\(\int d_s|\epsilon_\KK|\) and \(\int d_w|\epsilon_S|\).
The mass differences satisfy logistic comparisons with
coefficient at most three on \([0,2]\), plus
\(\int|\epsilon_\KK|\) and \(\int|\epsilon_S|\).
Thus the same coefficient calculation as CWI gives
\[
 D^+\dist\le(1+40R)\dist+\mathcal R
                \le64R\dist+\mathcal R.
\]
All differentiations are justified by the finite second moment
and the integrable cubic majorants proved above, or by truncation.
No fixed-relative-weight assumption is imposed on the original
flow. Gronwall, the entry distance, and
\eqref{an07fwcv1:reductionnorm} prove \eqref{an07fwcv1:distbound}.
The logit consequence also follows directly from Lemma
\ref{an07fwcv1:massfine}, with the identical initial phase.
\end{proof}

\section{Background precision and A-error-transfer}
\label{an07fwcv1:errorsection}
\begin{lemma}[Background on the continued window]
\label{an07fwcv1:background}
For \(T_0=\Sigma\cup\BB\), uniformly on \(I_q\) and
\(0\le\zeta\le R_\zeta\),
\begin{equation}
 \EE_{T_0}\le C_Dq^2\epsilon_{\rm bg},\qquad
 \mathcal Q_{T_0}\le C_{D,R_\zeta}q^2\epsilon_{\rm bg},\qquad
 |E_{\xi_q}-E_{\xi_q}^{\KK}|\le
                   C_{D,R_\zeta}q^2\epsilon_{\rm bg}.
 \label{an07fwcv1:backgroundbound}
\end{equation}
\end{lemma}
\begin{proof}
All four PSE \texttt{conditional} hypotheses have been supplied:
cross entries and \(R_{22}^{(1)}\) by
Lemma \ref{an07fwcv1:transverse}, alignment by
Lemma \ref{an07fwcv1:charts}, and the lower gate by loss and
the same Gaussian calculation as PSE:
\[
 |R_{22}^{(2)}|\le2e^{-169\sin^2\theta/(3200q^2)}
                                      \quad(U_2<0).
\]
Its proof uses only the exact gate probability, bounded residual
\(L^2\), and \(\E X_2^4\le43/16\), so it is valid after
the old passage. Apply the conditional inequality with
\(a=t_n\), duration \(D\), and its parameter \(w=1\).
The initial energy is in \eqref{an07fwcv1:entry}; the rough
mass comparison gives \(\mu_{T_0}\le C_Dq^{\gamma_*}\).
Hence the added energy is at most \(C_Dq^2Lq^{\gamma_*}\),
bounded by the displayed \(q^2\epsilon_{\rm bg}\).
The exact static inequality
\(\mathcal Q_{T_0}\le2q^2\zeta^2\mu_{T_0}+2\EE_{T_0}\)
gives the second bound, including every changing-angle slice.

Primary inactivity gives \(f=f_\KK+f_{T_0}\) at the probe.
CWI's exact split is
\[
 E-E^\KK=-\mathcal Q_{T_0}+\mathcal A_{T_0}
       +f_\KK\cdot f_{T_0}+\tfrac12|f_{T_0}|^2,
 \quad |\mathcal A_{T_0}|\le C_DqL\sqrt{\mu_{T_0}\mathcal Q_{T_0}},
 \quad |f_{T_0}|\le\sqrt{\mu_{T_0}\mathcal Q_{T_0}}.
\]
The bounded core chart gives \(f_\KK=O_D(q)\).
Young's inequality and \(L^2\mu_{T_0}=o(1)\) pay all
cross terms by \(C_Dq^2(q^8+L^2q^{\gamma_*})\).
No favorable sign is substituted for an absolute precision bound.
\end{proof}

\begin{lemma}[Gap-free core static comparison]
\label{an07fwcv1:static}
For \(R\ge2\), strong mass \(M\le2\), strong coordinates
and the comparison atom in \([-R,R]^2\), and \(|\zeta|\le R\),
put
\[
 A=\int_\KK(\zeta-x)_+\dd\mu,\quad
 B=\int_\KK y(\zeta-x)_+\dd\mu,\quad
 \mathcal E_\KK=-\zeta A+B+\tfrac12A^2.
\]
Then
\begin{equation}
 E_{\xi_q}^{\KK}-\tfrac12=q^2\mathcal E_\KK+O_R(q^4),
 \qquad |\mathcal E_\KK-\mathcal E_*|\le10R^2\dist.
 \label{an07fwcv1:staticcompare}
\end{equation}
No condition on the sign or size of \(\zeta-X\) is required.
\end{lemma}
\begin{proof}
The input activation is \(\sin(q(\zeta-x))_+\).
Taylor expansion with the 1-Lipschitz positive part gives
\((f_\KK)_1=qA+O_R(q^3)\),
\((f_\KK)_2=q^2B+O_R(q^4)\). Expand the exact squared error.
For arbitrary gap, take \(A_*=(\zeta-X)_+\),
\(B_*=Y_*A_*\), not \(\zeta-X\).
Then
\[
 |A-A_*|\le D_s+2R|M-1|,\quad
 |B-B_*|\le2RD_s+2R^2|M-1|,
 \quad |A|\le4R,\ |A_*|\le2R.
\]
Consequently the energy difference is at most
\(4R|A-A_*|+|B-B_*|\le10R^2\dist\).
This is the static part of CWI's proof with its unnecessary
rate-gap hypothesis removed by the displayed argument.
It covers atoms exactly at the probe gate.
\end{proof}

\begin{theorem}[Finite-window continuation and full A-error-transfer]
\label{an07fwcv1:main}
For the original initialized deep-regime family, retain the proved
entry level and center \eqref{an07fwcv1:time}. For every fixed
\(T>\max_\rho s_n\) and \(R_\zeta<\infty\), the primitive
continuation conclusions \eqref{an07fwcv1:primitive}--
\eqref{an07fwcv1:masslogistic}, the field and weighted reaction
reduction \eqref{an07fwcv1:reductionnorm}, and
\eqref{an07fwcv1:distbound} hold on \([t_n,\tau_h+T]\).
There are no additional retained original-flow hypotheses.
Moreover
\begin{equation}
 \sup_{\substack{\rho\in[13/20,7/10],\ s_n(\rho)\le s\le T\\
                       0\le\zeta\le R_\zeta}}
 \left|q^{-2}(E_{\xi_q}(\tau_h+s)-\tfrac12)
              -\mathcal E_{\rho,\zeta}(s)\right|
 \le C_{T,n,R_\zeta}q^{1/20}\longrightarrow0.
 \label{an07fwcv1:errortransfer}
\end{equation}
Together with Corollary \ref{an07fwcv1:early}, this gives H2's
A-error-transfer on every fixed finite centered window,
uniformly on the stated ratio box and every bounded nonnegative
offset interval. It supplies LP's \texttt{compact-transfer}
premise on such compact offsets, with
\(t_h^{\rm phys}=\tau_h/\mu_1^2\).
\end{theorem}
\begin{proof}
Combine Lemmas \ref{an07fwcv1:background} and
\ref{an07fwcv1:static} with Theorem \ref{an07fwcv1:propagation}.
The error is at most
\(10R^2\dist+C_Rq^2+C_{D,R_\zeta}\epsilon_{\rm bg}\).
Since \(\gamma_*>1/20\),
\(q^{-1/20}\epsilon_{\rm bg}\to0\), proving the stated rate.
For a general fixed window, the part \(s\le s_n(\rho)\)
is covered by the same compact-level argument as the early
corollary (its minimum level is uniform); the remaining part
is covered by \eqref{an07fwcv1:errortransfer}.
The center is never replaced by \(t_b\) or by an order-one
approximation to a seed. The known difference \(Q-\lambda=q^2\)
costs \(O(q^2\log(1/q))=o(1)\) in the inherited long entry phase;
the new interval compares the actual mass with the exact same
\(\lambda\)-logistic. Normalized time in H2 and LP's physical
time agree since
\(\tau_h+s=\mu_1^2(t_h^{\rm phys}+s/\mu_1^2)\).
No later-window rate or growing-window conclusion is used.
\end{proof}

\section{S5d handoff and final constant audit}
\label{an07fwcv1:handoff}
The new constant order, appended to the unchanged WPT/RB--WCS
order, is
\[
 (T,R_\zeta,D)\ \longrightarrow\ \bar S\
 \longrightarrow\ C_Z,C_X,C_\Delta,R_A,A_D\
 \longrightarrow\ R\ \longrightarrow\ C_{\rm red},C_{\rm bg}\
 \longrightarrow\ q.
\]
Here \(R\) covers the leading solution on the fixed window as
well as the original core and offsets. Each primitive guard is
chosen after a bound independent of itself. Neither an old small
energy guard nor an order-one passage allowance is hidden in a
vanishing pointwise defect. All fixed-time losses may depend on
\(D\); no uniformity as \(T\to\infty\) is claimed.

\small
\begin{longtable}{>{\raggedright\arraybackslash}p{.27\textwidth}
 >{\raggedright\arraybackslash}p{.63\textwidth}}
\toprule S5d consumer & Supplied quantity and remaining action\\\midrule\endhead
Same original state and phase
& Physical masses and current first moments obey
 \(\dist=O_T(q^{1/20})\); original selected mass/logit uses the
 unchanged \(\tau_h\). No passage bound is being used after \(t_b\).\\
Core sourcewise field/reaction
& Equation \eqref{an07fwcv1:corefield} has pointwise
 source errors \(O_T(q^{1/10})\); source reactions are given
 there and in Lemma \ref{an07fwcv1:massfine}. All residual
 producers, including probe-inactive weak labels, remain included.\\
Leading rate robustness
& CWI \texttt{rate} applies once the learned parameter/window
 and gap hypotheses are matched: positive measures, masses
 \(\le2\), strong chart and comparator in \([-R,R]^2\), finite
 weak first moments, \(|\zeta|\le R\),
 \(\zeta-X\ge g_0=g_*/4>0\). It supplies
 \(64R^2(1+g_0^{-1})\dist\). The learned-window and margin
 choices themselves belong to WIT/S5d.\\
Background source charges
& PSE/CWI matrix and alignment hypotheses now hold on the
 full window. CWI \texttt{ratecharge} therefore gives
 \(|\xi_q\cdot\dot f_{\Sigma\cup\BB,p}|
 \le C_Tq^2\epsilon_{\rm bg}\), and the associated unprojected
 bound. Primary direct outputs/rates vanish by strict inactivity.
 These are source charges, not the asserted A-rate-transfer theorem.\\
Gate conventions
& Static transfer needs no gap. For S5d's uniform source transfer,
 use WIT's gate-separated \(J\) and its \(\mathcal P_*\);
 verify \(g_0\), source and total-rate margins in that range.
 Retain selector/a.e. statements elsewhere. No pointwise
 gate-hit-atom rate claim is added.\\
H2 and LP
& A-error-transfer and LP's finite-window premise are supplied
 here for the initialized family and matched physical center.
 Applying WIT/LPR/LP to obtain the final reversal/persistence
 statements and AN8 assembly remains outside this increment.\\
\bottomrule
\end{longtable}\normalsize
The analytic estimates needed for S5d's original-flow rate calculation
are supplied by the sourcewise reduction and first-moment comparison;
the next step must assemble them on the precise learned windows and
compare with their margins. This increment does not assert that
rate-transfer conclusion. All proofs above use Gaussian identities,
differential inequalities and fixed-time comparison only.
\end{document}
